\subsubsection*{Representing DPFs} \label{subsub:implementationRepresentingDPFs}
Both, Distributed Point Functions (\ref{sub:DistributedPointFunctions}) and Distributed Comparison Functions (\ref{sub:DistributedComparisonFunctions}), inherit from the same parent class \code{dpf}.
Again, recall the representation of point- and comparison-functions as binary trees.
Thus, the most trivial representation of a single point function is saving the seed of the root node and generating the whole tree every time the point function is used.
Since this approach is too computational expensive for repetitive evaluations, DPFs are stored in an already evaluated form with the possibility to recreate the tree when needed.

Therefore, the base class \code{dpf} stores the seed of the root node.
Also, it stores a vector \code{cfbits} which holds all flag bits on the path to target index $\alpha$.
Bit $i$ of this vector represents the bit of the $i$'th node on the path to $\alpha$.
It also implements the functionality to descend a tree given a seed and a bool indicating which child to calculate.
The code further defines special DPFs namely \textit{Distributed Point Functions for Memory Access (RDPF)} and \textit{Distributed Comparison Functions (CDPF)}.

\paragraph{RDPFs}\label{par:preliminariesRDPFs}
RDPFs inherit from the basic DPF implementation.
It expands its parent with further attributes and the \code{LeafInfo} struct.
\code{leaf\_cfbits} stores the correction-flag bits for the leaf level.
Recall the correction-flag bits being the low words of the labels (\ref{sub:BGIDPFConstruction}).
This is equivalent to $t_b$ from the high-level \textit{DUORAM} overview in \ref{sec:3pDuoram}, which means bit $i$ holds the correction flag-bit of leaf $i$.
\code{li} is a vector holding \code{LeafInfo} structs.
In normal RDPFs, which is the standard in this research, it will only hold a single \code{LeafInfo} struct.
\code{LeafInfo} is essential for using the evaluated RDPF result in a read or write operation.
It holds the correction word for the leaf level, which itself is a XOR-share, and all further attributes to easily convert it into an additive share.
This is essential, since \code{duoram} is based on additive shares. 
Thus converting from XOR RDPF-shares to additive shares simplifies further evaluation processes a lot.
More information on \code{LeafInfo} and the concept of \code{DPFNode}s and \code{LeafNode}s is given in \ref{subsub:DpfNodeAndLeafLabel}.

The code stores RDPFs as triples, called \code{RDPFTriple}.
As the name suggests, it holds three RDPFs, each with the same target, again, either XOR or additively shared.
To generate a RDPF-Triple including three RDPFs on a given target, one uses its constructor which itself calls the RDPF's constructor three times.
All three RDPFs are stored in a vector, the \code{RDPFTriple}.
The triple itself is stored in the \code{mpcio} class assigned to each player.
When generating these triples in the preprocessing phase, they are written to the background memory and called automatically when performing reads or updates in the online phase.
This construction makes it easy to keep track of the used RDPFs.

\paragraph*{Leaf Construction}\label{par:leafConstruction}
Leafs are stored as concatenations of 64-bit vectors.
A single leaf holds the functionality for both reading and writing.
This means it holds 64-bit for writing, which will add up to $M_b$ at the target index, and 64-bit for reading, adding up to the shifting vector when combined with the other peer's shifting vector.
Concatenating both 64-bit values in a single value results in a single 128-bit vector.
Writing bits are placed at the upper 64-bits, the lower word represents shares of the shifting vector.
Therefore, when executing a read and requiring the flag bits, one can simply add up (or XOR) the least significant bits of all the leafs without performing additional shifting operations.
This means a tree of depth $d$ holds $2^d$ leafs, each of size 128.

\paragraph{CDPFs}\label{par:preliminariesCDPFs}
CDPFs too, inherit from the base structures predefined by \code{dpf}.
They are, however, stored as tuples.
One can choose whether to share the target additively via \code{as\_target} or distribute XOR-shares, called \code{xs\_target}.
This process is similar to RDPFs, since both CDPFs within a tuple hold the same target.
Thus, combining the targets of $P_0,P_1$ results in the actual target index.
A tuple of CDPFs allows the client to compare a given value to $0$ as described in \ref{sub:CDPFConstruction}.

\subsubsection*{types.hpp} \label{subsub:implementationTypeshpp}
This file defines basic macros, aliases, structs, and their operators.

The most prominent type-aliases are:
\begin{itemize}\label{itemize:importantAliases}
    \item \code{nbits\_t}: 
    This value counts the number of bits in a value, which means \code{value\_t}.
    Thus, the chosen datatype needs to be large enough to store \textit{VALUE\_BITS} bits.
    The default value is \code{uint8\_t}.
    \item \code{bit\_t}:
    Chooses how a single bit is represented, usually in the \code{RegBS} struct.
    The default datatype for this alias is \code{bool}.
    \item \code{address\_t}: 
    Decides the size of the addresses in the secret-shared memory.
    Is set to \code{uint32\_t}.
    \item \code{DPFnode}: 
    Implements the datatype of a DPFnode and establishes the basis for defining Leaf nodes.
    It is of type \code{\_\_m128i}.
    Leafs are also implemented using this type.
    Evaluation on leaf nodes splits the vector in a lower and upper half which makes it two 64-bit vectors.
    More information on this alias is provided in \ref{subsub:DpfNodeAndLeafLabel}.
    \item \code{value\_t}: 
    Decides the datatype of values in the secret-shared memory.
    The default value is \code{uint64\_t}.
    \item \code{AES\_KEY}:
    Represents an AES-key for 128-bit AES.
    More details on this alias are found in \ref{subsub:implementationAesni} and \ref{sub:AESSizeLimitations}.
\end{itemize}
Important macros are:
\begin{itemize}\label{itemize:importantMacros}
    \item {VALUE\_BITS}: 
    This will decide the size of \code{value\_t}. 
    The default value is 64.
    \item {ADDRESS\_MAX\_BITS}: 
    Decides the size of \code{address\_t}.
    The default value is 32.
\end{itemize}
The most prominent structs are:
\begin{itemize}\label{itemize:importantStructs}
    \item \code{RegAS}: 
    Represents a register holding an additive share of type \code{value\_t}.
    \item \code{RegXS}: 
    Represents a register holding an XOR share of type \code{value\_t}.
    \item \code{RegBS}: 
    Represents a register holding a bit-share of type \code{bit\_t} which itself is a \code{bool} type-alias.
\end{itemize}

\subsubsection*{Intel AES-NI} \label{subsub:implementationAesni}
The existing source code is optimized for Intel's \textit{Advanced Encryption Standard New Instruction Set (AES-NI)}\footnote{Intel's publication on AES-NI is published on their \hyperlink{https://www.intel.com/content/dam/doc/white-paper/advanced-encryption-standard-new-instructions-set-paper.pdf}{website} (https://www.intel.com/content/dam/doc/white-paper/advanced-encryption-standard-new-instructions-set-paper.pdf)}.
As the name suggests, it extends and optimizes Intel's existing instruction set for AES-operations.
These improvements come in form of the \code{wmmintrin.h} and \code{emmintrin.h} header files implementing core functions of \textit{AES-ECB (Advanced Encryption Standard in Electronic Code Book Mode)} based on \textit{AES-NI}.
As this header file is included in the gcc compiler, no additional installations for using both headers are required.
Yet, the flags \code{-maes}, and \code{-msse2} need to be included for the linker to work.
Since AES-ECB is of great use when generating DPFs, these optimizations will foster great performance improvements.
Operations like \code{AES\_key\_gen\_assist} are optimized for the \code{\_\_m128i} \textit{Single Instruction Multiple Data (SIMD)} datatype, representing a 128-bit vector.
Therefore, the implementation stores any data used in AES-contexts in this format.

The careful reader notices \code{\_\_m128i} vectors being twice as large as the \code{uint64\_t} type-alias \code{value\_t} (\ref{itemize:importantAliases}).
This is due to AES-ECB being used as a \textit{Pseudo Random Generator (PRG)} during creation and traversion of DPF trees.
The code uses its security properties of being computationally random \cite{rijmen2001advanced} to calculate new labels given a seed and key.
Both parties use the same private key, thus AES-ECB becomes pseudorandom.
Labels are stored as \code{\_\_m128i}, leafs as explained in the previous section \ref{par:leafConstruction}.
Therefore, by using a \code{\_\_m128i} seed and labels encrypting each label using AES-ECB recursively, one simulates a length doubling pseudo random generator \cite{vadapalli2023duoram} as required to run the BGI-DPF-Construction (\ref{sub:BGIDPFConstruction}).
More on label and leaf construction can be found in the following sections.

AES for 128-bit uses ten encryption-rounds \cite{rijmen2001advanced}.
Therefore, after inserting the seed which resembles the key of the AES operation, ten more sub-keys are generated starting from the seed.
This structure is stored in format of an \code{\_\_m128i[11]} array.
The seed at index $0$, the ten sub-keys at the following indices.
It is represented as type-alias \code{AES\_KEY} which is equivalent to \code{\_\_m128i[11]} (\ref{itemize:importantAliases}).
To simplify key-creation, \code{AES\_KEY} is wrapped in the struct \code{PRGKey}.
It takes a single \code{\_\_m128i} value as constructor input and generates ten more sub-keys using the input as seed for AES-ECB.
As the input and the key are set to the same value every time, one receives a pseudorandom number generator.

For AES to be pseudorandom, keys need to be the same throughout instances.
Thus, all instances use the same \code{PRGKey} instances for creating labels and leafs, called \code{prgkey} and \code{leafprgkeys[3]}.
Note \code{leafprgkeys[3]} being an array holding three \code{PRGKey} structs.
This distinction between keys for label and leaf generation is necessary as one needs to differentiate cases handling different values for \code{LWIDTH} (\ref{subsub:WIDTHandLWIDTH}).
A non-standard value of \code{LWIDTH}, which means any value $x\neq1$, will require multiple keys in the leaf layer as several Leaf labels need to be instantiated (\ref{subsub:DpfNodeAndLeafLabel}).
As this matter is only present in the leaf layer, using separate keys during generation will increase performance and readability of the code.
The implementation uses the first 20 digits of $\pi$ for \code{prgkey} and 20 digits of $e$ for each instance of \code{leafprgkeys}.

Requirements derived from this, along with further details on \code{WIDTH} and \code{LWIDTH} are described in \ref{sec:GeneralLimitations}.

One practical drawback of using \code{\_\_m128i} vectors is its debugging incompatibility when using \textit{Valgrind}\footnote{More information can be found on their website and documentation at https://valgrind.org/}.
\textit{Valgrind} is a debugging suite designed for spotting memory leakages in the code.
It is, however, not compatible for using SIMD datatypes which triggers an error and ultimately stops program execution.

\subsubsection*{Duoram and Shapes} \label{subsub:implementationDuoram}
Each party is assigned a \code{duoram} object.
It holds the following information: 
\begin{itemize}
    \item player: 
    The number of the party.
    0 or 1 for computational peers and 2 for the server.
    \item database: 
    Vector of registers representing the party's share of the database $D_b$.
    These entries are either additively (\code{RegAS}) or XOR shared (\code{RegXS}).
    Entries are stored as \code{RegAS} by default and can be changed manually.
    This means the database is a vector holding $n$ register entries.
    \item blind: 
    The blinding factor $\zeta_b$ to blind its own database.
    As the blinding factor is continuous, it has the same size as the database.
    \item {peer\_blinded\_db}:
    This is the blinded database of party $P_b\text{ with }b\in\{0,1\}$ the other party $P_{b'}\text{ with }{b'-1}\in\{0,1\}$.
    \item Party $P_2$ uses the attributes \code{p0\_blind}, \code{p1\_blind} to represent the blinded databases of both parties $P_0,P_1$.
    The size is, again, the same as the database.
\end{itemize}

Each \code{duoram} must be accessed via a hierarchy of \textit{Shapes}.
The implementation uses a hierarchy consisting of solely a single shape called \code{Flat}.
A shape provides basic functions for reading and updating the duoram.
This is primarily the index operator \code{[]} for reading and the \code{+=} operator for updating.
This has the great advantage of updating accessing contexts while keeping the underlying database consistent.
For example, when executing multiple database accesses multiple RDPF-Triples and thus separate contexts are needed as keeping the context might leak accessing-information due to reusing the same RDPF-Triple.
Instead of assigning these contexts to the same database introducing a great overhead one simply updates the underlying shape introducing a fresh context.
Without a shape, one manually has to update the database, reassign new RDPF-Triples to each player and manipulate their \code{mpcio} objects.
\code{context()} as defined in \code{Flat} generates a new shape including fresh RDPF-Triples and assigns it to each party.
This reduces the overhead to a minimum and introduces independent reads and writes.

The base implementation uses the \code{Flat} shape only.
Extensions on different shapes are not planned or necessary, thus no further descriptions on different shapes and shape hierarchies are required.

\subsubsection*{WIDTH and LWIDTH}\label{subsub:WIDTHandLWIDTH}
RDPFs can be reused for reading.
This, however, will leak the information that reads were executed on the same location multiple times.
The location stays secret \cite{vadapalli2023duoram}.

This process is similar for writing, yet it requires a few code adjustments.
In short, one simply extends the given leaf construction (\ref{par:leafConstruction}) by adding as many 64-bit values representing shares of $M$ as wanted.
Vadapalli et al. chose to template their \code{RDPF} struct with a property called \code{WIDTH} of type \code{nbits\_t} (\ref{itemize:importantAliases}).
\code{WIDTH} is set to $1$ by default.
This indicates one write access per RDPF, no reuse is possible.
Increasing this value to $n$ adjusts the type aliases \code{RegASW} and \code{RegXSW}, both containing $n$ instances of \code{RegAS} and \code{RegXS} respectively.
On a high level, this means that each RDPF can hold up to $n$ different values for the message $M$.
Therefore, by increasing \code{WIDTH} to a value of $n$, one can reuse a given RDPF for up to $n$ different writing accesses.

\code{LWIDTH} is directly derived from \code{WIDTH}.
It is calculated by $1+\text{WIDTH}/2$.
As it is of type \code{nbits\_t} (\ref{itemize:importantAliases}) setting \code{WIDTH} to $1$ will set \code{LWIDTH} to $1$ too.
Recall leaf nodes being stored as an array of \code{\_\_m128i} values, with 64-bit reserved for reading and 64-bit reserved for writing.
Performing either a single read or write will only require the array to contain only a single \code{\_\_m128i} value.
Increasing \code{WIDTH} to $n$ will lead to $n$ 64-bit values being needed for $n$ writes.
Therefore, the array representing leaf nodes must be adjusted to hold $n$ \code{\_\_m128i} values.
\code{LWIDTH} represents the number of \code{\_\_m128i} values necessary to accommodate \code{WIDTH} many writes per instance.

Limitations on these values are listed in \ref{sub:LimitWIDTHandLWIDTH}, more information on this construction is given in the following section.

\subsubsection*{DPFNode and LeafNode}\label{subsub:DpfNodeAndLeafLabel}
The concepts of DPF nodes, leafs, \code{WIDTH}, and \code{\_\_m128i} vectors are closely intertwined.

As already mentioned, DPF trees are generated using AES-ECB (\ref{subsub:implementationAesni}).
For the goal of storing 64-bit values in the base implementation, one needs to simulate a length-doubling PRG.
This can be done by using AES-ECB on vectors twice as large as the actual label, which means 128-bits.
Therefore, a single node for a particular DPF tree storing values of size 64 has nodes of size 128.
Vadapalli et al. consequently drew the logical conclusion to use \code{DPFnode} as type alias for \code{\_\_m128i}.
As Intel's AES-ECB implementation (\ref{subsub:implementationAesni}) directly uses \code{\_\_m128i} vectors as input and output, \code{DPFnode} can be used for AES operations without conversion.

Recall the definition and usage of \code{WIDTH} and \code{LWIDTH} from \ref{subsub:WIDTHandLWIDTH}.
In short, each leaf can be divided into $2n*$\code{LWDITH} blocks of size 64, which makes it reusable for \code{WIDTH} many updates.
This means, to implement a DPF tree with a certain value of \code{WIDTH}, one needs a leaf label of size $\code{LWIDTH}*128$.
Thus, a \code{LeafNode} is chosen to be a type alias for a vector holding \code{LWIDTH} many \code{\_\_m128i} values or \code{DPFnode}s respectively.
This results in the type definition \code{using LeafNode = std::vector<DPFNode, LWIDTH>}. 

\subsubsection*{mpcio and Communication} \label{subsub:implementationMpcioAndCommunication}
All three parties execute the same executable \code{prac} which can be generated by using the \code{MAKEFILE}.
These executables communicate asynchronously via TCP-sockets.
When executed locally, which this research assumes to be the default, the peers addresses are all hosted on the \code{localhost}.
Sockets, player information, sent and received messages and the functionality to send and receive are implemented in the \code{MPCPeerIO} struct.
\code{MPCPeerIO} can be freely translated to \textit{"Input and Output Functionality for performing MCP communication with computational peers"}.
It is one of the first structs to instantiate and will be used to create most of the other objects.

\code{MPCPeerIO} is the child of the \code{mpcio} struct which holds general io-attributes and also provides functionality for timing the communication process and timing of round-trip-trimes.
This struct also holds information about the chosen processing mode.
The program offers three possibilities for this value: 
\begin{itemize}
    \item -p: 
    Defines the \code{PRECOMPUTATION\_MODE}, which will execute the preprocessing phase as explained in section \ref{subsub:DuoramPreprocessing}.
    \item leaving the specifier blank: 
    Will execute the \code{ONLINE\_MODE}.
    Note that this requires the preprocessing mode to be executed first since it requires dpf-triples for reads and writes.
    Thes user can specify which tests to execute.
    The default test is called \code{duoram} and will check reads and writes, each independant and interleaved.
    Read more about testing in section \ref{sub:testingAndBenchmarking}.
    \item -o:
    Will execute the \code{ONLINE\_ONLY} mode.
    This means no precomputation will take place, all needed resources are computated on the fly.
    Of course, this will impact performance but will come in handy for performing smaller tests.
    With this specifier, too, one can choose which tests to execute.
    These tests will be the same as in the online mode.
\end{itemize}

All three parties will use a predefined message structure.
Therefore, generated precomputation \ref{subsub:implementationPreprocessing} data will implement this protocol.
It separates a message into byte-segments with each segment accomodating a special type of data:
\begin{itemize}\label{itemize:preprocessingStructure}
    \item 0x01: Indicates the type.
    \item 0x01 to 0x30: RDPFs. Since the size of RDPFs depends on its depth the given space is defined in this manner.
    \item 0x40: CDPFs 
    \item 0x80: Multiplication Triple for performing MPC-multiplications (\ref{sub:MultiplicationAdditionTriples})
    \item 0x81: AND Triple for performing MPC-AND-operations (\ref{sub:MultiplicationAdditionTriples})
    \item 0x83: Select Triple. Since it will not be uesd further the benevolent reader may read the comments in \code{types.hpp}.
    \item 0x8e: Counter counting the AES-operations. This can be used for testing.
    \item 0x00: Shows the end of preprocessing. 
    \item One extra byte is reserved for \code{subtype}.
    More details about this byte can be found in \ref{sub:LimitWIDTHandLWIDTH}.
\end{itemize}

More detailed information about the TCP-sockets and the underlying functions is given in \ref{subsub:gmpCommunication}.